\section{Cicero 的结果与闭环架构}

在理解任务后，才能正确解读 Cicero 的成绩。系统使用约 50,000 局人类游戏数据建立语言与行为模型，并在在线联赛匿名参加 40 局。课堂刻意同时展示定性案例和定量表格，因为“像人”与“打得强”是两条不同证据。

\subsection{系统定位：会沟通的 Diplomacy agent}

slide 21 给出最简定义：Cicero 把战略推理与自然语言生成结合，用于在线 high-level play。这里的“AI agent”不是只有一个模型 endpoint，而是持续接收 board、orders 和 messages，再产生行动与消息的循环系统。

\lecturefigure{slide-21-cicero-intro.jpg}{Cicero：结合 strategic reasoning 与 grounded dialogue 的 Diplomacy agent}{官方录像恢复幻灯片 V021；Bakhtin et al., Science 2022}

\subsection{定性证据：玩家是否把它当作参与者}

系统定义只说明 Cicero 打算做什么，不能证明模块在真实联赛中形成了可信行为。本节先看定性证据，再与下一节的分数表交叉验证。第二十二张 slide 展示匿名参赛和玩家反馈。一个好的 qualitative result 不是“消息语法正确”，而是其他玩家愿意依据消息协调行动，并在长局中持续把 Cicero 当作有策略、可谈判的参与者。课堂也保留边界：有人赛后表达过怀疑，但并未形成系统性识别。

\lecturefigure{slide-22-cicero-human-qualitative.jpg}{Cicero 与人类在线对局的定性反馈}{官方录像恢复幻灯片 V022；Stanford CS25 Lecture 14}

\subsection{定量证据：强，但不宣称 superhuman}

定性反馈说明系统没有频繁破坏协作，但强度必须回到可比较的数值。本节读表时要同时保留样本量、参与者门槛和排名口径。slide 23 给出更具体的审计数据：40 局、82 位 unique players、每局平均收发 292 条消息；在至少参加五局的 19 位玩家中，Cicero 平均得分 25.8\%，排名第二，并进入全部参与者前 10\%。这些数字证明其在线竞争力，但样本、联赛构成和评分规则仍不足以支持无条件“超人”结论。

\lecturefigure{slide-23-cicero-human-results.jpg}{Cicero 在线联赛结果：25.8\% 平均分、40 局与前 10\% 排名}{官方录像恢复幻灯片 V023；Stanford CS25 Lecture 14}

\begin{warningbox}{不要把“未被识别”写成通过图灵测试}
实验是在真实联赛中匿名对局，不是预注册的 bot-detection study。结果能说明消息没有频繁暴露明显机器模式，却不能证明人类无法识别 AI，也不能单独证明语言理解达到人类水平。
\end{warningbox}

\teachervoice{老师主动使用“strong human-level”而不是“superhuman”。这是一条重要的证据纪律：系统结果足够令人兴奋时，更应说明比较群体、局数和评价条件，而不是用更强形容词替代实验设计。}

\subsection{闭环架构：预测、规划、意图、生成、过滤}

架构图从左到右并不是一次性 pipeline，而是每收发一条消息都会重跑的 recurrent loop。board state、action history 和 dialogue history 进入 action prediction；human imitation policy 经 piKL search 调整；规划出的近程行动形成 intents；dialogue model 生成候选消息；过滤器剔除 nonsense、矛盾和战略上有害的候选。

\lecturefigure{slide-24-cicero-architecture.jpg}{Cicero 完整闭环：从状态与对话到行动预测、规划、意图和消息过滤}{官方录像恢复幻灯片 V024；Bakhtin et al., Science 2022}

若 $s$ 是 board state，$h$ 是 action history，$d$ 是 dialogue history，七位玩家下一回合行动的联合预测可以写成
\begin{equation}
p(a_1,\ldots,a_7\mid s,h,d).
\end{equation}
这里 $a_i$ 往往是某位玩家多个单位 orders 的组合。预测不是最终目标，它既给 search 提供他人策略模型，也决定哪些消息可能真正改变联合行动。

\begin{importantbox}{模块化不等于模块彼此独立}
战略模型与语言模型职责分开，但通过 intent 和 predicted response 紧密耦合。若完全端到端，语言模型同时承担世界建模、策略选择和表面表达，难以控制；若完全分离，语言又无法改变计划。Cicero 的设计是在中间接口处交换结构化意图和行为分布。
\end{importantbox}

\subsection{本章小结}

Cicero 的在线结果显示系统能在高消息密度的人类联赛中稳定工作，但课堂没有把它夸大为无条件超人。架构的关键是闭环：语言不是最终包装，而是会改变其他玩家 policy、随后被重新纳入规划的行动。

